\documentclass[11pt]{amsart}
\usepackage[margin=1in]{geometry}
\usepackage{amsmath,amssymb,amsthm}
\usepackage{microtype}
\usepackage[hidelinks]{hyperref}
\newtheorem{theorem}{Theorem}
\newtheorem{lemma}[theorem]{Lemma}
\newtheorem{proposition}[theorem]{Proposition}
\newcommand{\VR}{\operatorname{VR}}
\newcommand{\ZZ}{\mathbb Z}
\newcommand{\st}{\operatorname{st}}
\newcommand{\conv}{\operatorname{conv}}
\title{Torus-grid Vietoris--Rips complexes at half the diameter}
\author{Kaiwen Jin}
\email{yui.ui.siki@gmail.com}
\date{October 14, 2026}
\hypersetup{pdftitle={Torus-grid Vietoris-Rips complexes at half the diameter},pdfauthor={Kaiwen Jin}}
\begin{document}
\begin{abstract}
For every integer \(k\geq3\), we determine the homotopy type of the closed
Vietoris--Rips complex of the \(2k\) by \(2k\) periodic Manhattan grid at
radius \(k\). It is
\(S^3\vee\bigvee^{4k}S^{2k-1}\), as conjectured by Adams, Adetowubo,
Barriga-Acosta, Feng and Sterling. Adjacent pairs of rows and columns give
the high-dimensional spheres. Removing one facet from each leaves a
complex with the integral homology of a three-sphere. The comparison is
proved uniformly through a homotopy calculation for finite cylinders
and a cover by coordinate semicircles.
\end{abstract}
\maketitle

\section{Introduction}

Vietoris--Rips complexes of periodic samples provide finite models in which
one can follow the change from local geometry to higher-dimensional
topology. For a torus grid with the periodic Manhattan metric, Adams,
Adetowubo, Barriga-Acosta, Feng and Sterling~\cite{torus} determined several
families of homotopy types. Their Theorem~5.8 identifies the small-scale
torus regime, while Theorems~5.10 and~5.12 describe the first two
transition families. Their Question~8.6 asks for the homotopy type at
half the diameter of an even square grid. We answer that question for all
\(k\geq3\); the case \(k=2\) was already known.

For \(n\geq1\), write
\[
T_{n,n}=(\ZZ/n\ZZ)^2,\qquad
\delta_n(a,c)=\min\{|a-c|,n-|a-c|\},
\]
using representatives \(0,\ldots,n-1\), and set
\[
d((a,b),(c,e))=\delta_n(a,c)+\delta_n(b,e).
\]
The closed complex \(\VR(T_{n,n};r)\) contains a finite subset precisely
when every pair in it has distance at most \(r\). It is the clique complex
of its graph of edges.

\begin{theorem}\label{thm:main}
For every integer \(k\geq3\),
\[
\VR(T_{2k,2k};k)\simeq
S^3\vee\bigvee_{j=1}^{4k}S^{2k-1}.
\]
\end{theorem}

The metric diameter here is \(2k\), so the radius is exactly half the
diameter. The wedge is taken at a common base point. For \(k=2\) the same
formula gives nine three-spheres, agreeing with the known homotopy type
in~\cite{torus}.

The \(4k\) high-dimensional spheres have concrete representatives:
each pair of adjacent rows or columns induces a cross-polytope boundary.
We remove one suitable facet from each of these spheres. Their remaining
union has a good cover whose nerve is the join of two circles. To compare
this union with the rest of the Rips complex, we use coordinate
semicircles. The main calculation is a cylinder whose circular direction
has circumference \(2k\); adding a column adds exactly one
\((2k-1)\)-sphere. This also explains why the equality at the closed
threshold matters.

Throughout the proof, fix \(k\geq3\), put \(n=2k\), and write
\(K=\VR(T_{n,n};k)\). All homology groups have integer coefficients,
and simplicial complexes are identified with their geometric realizations
when discussing homotopy.

\section{Adjacent rows and columns}

For \(b\in\ZZ/n\ZZ\), let \(B_b\) be the induced subcomplex of \(K\) on the
two rows with second coordinates \(b,b+1\). For \(a\in\ZZ/n\ZZ\), let
\(D_a\) be the induced subcomplex on columns \(a,a+1\).
All indices in these definitions are taken modulo \(n\).

\begin{lemma}\label{lem:bands}
Each \(B_b\) and \(D_a\) is the boundary of an \(n\)-dimensional
cross-polytope. In particular, it is homeomorphic to \(S^{n-1}\).
The faces
\[
\tau_b=\{(a,b):a\ \text{even}\}
       \cup\{(a+k,b+1):a\ \text{odd}\},
\]
and their coordinate-swapped counterparts \(\upsilon_a\) in \(D_a\),
are maximal faces of \(K\). Each belongs to only its own band among the
\(2n\) complexes \(B_b,D_a\).
\end{lemma}

\begin{proof}
Every pair within one row is adjacent. Between the two rows the only
nonedges are
\[
(a,b),\quad (a+k,b+1).
\]
These form \(n\) opposite pairs, so the clique complex is the
cross-polytope boundary. The selected face chooses one vertex from each
pair.

To check maximality in \(K\), consider a vertex \((x,y)\) outside the
two rows. Its circular distance from at least one of them is at least
two. The first coordinates selected in either row form a parity class;
their maximum circular distance from \(x\) is at least \(k-1\).
Thus a vertex of the selected face is at distance at least \(k+1\)
from \((x,y)\). Inside the band, maximality follows from the opposite-pair
description.

The two occupied rows identify the horizontal band uniquely. Its
first-coordinate projection contains a parity class of \(k\geq3\)
points, so it cannot be contained in two adjacent columns. The vertical
claims follow by interchanging the coordinates.
\end{proof}

Let
\[
L=\bigcup_b B_b\ \cup\ \bigcup_a D_a.
\]
Delete the \(2n\) maximal faces \(\tau_b,\upsilon_a\), retaining all their
proper faces, from \(K\) and \(L\). Denote the resulting subcomplexes by
\(K^-\) and \(L^-\), respectively.

\begin{lemma}\label{lem:split}
There is a homotopy equivalence
\[
K\simeq K^-\vee\bigvee^{2n}S^{n-1}.
\]
\end{lemma}

\begin{proof}
In each band sphere, the complement of the interior of its selected facet
is a closed ball. By Lemma~\ref{lem:bands}, this ball contains no other
deleted facet. It therefore lies in \(K^-\), and contracts the attaching
map of the selected facet's boundary.
Restoring the facets one at a time gives the asserted wedge.
Equivalently, apply the null-homotopic union lemma
from~\cite[Lemma~5.7]{torus}: if both inclusions from \(X\cap Y\)
into \(X,Y\) are null-homotopic, then
\[
X\cup Y\simeq X\vee Y\vee\Sigma(X\cap Y).
\]
Here \(Y\) is the restored simplex and \(X\cap Y\) its boundary.
The complexes remain connected because their one-skeletons are unchanged.
\end{proof}

\begin{proposition}\label{prop:band-nerve}
The complex \(L^-\) is homotopy equivalent to \(S^3\).
\end{proposition}

\begin{proof}
Cover \(L^-\) by the \(2n\) band spheres with their selected facets removed.
Each member is a closed ball. Two horizontal bands intersect when their
indices are adjacent, in a full row simplex; three distinct horizontal
bands have empty intersection. The same holds vertically.
A mixed intersection is a simplex on a subset of a two-by-two grid,
whose pair distances are at most two. None of these intersections
contains a deleted facet.

Every nonempty intersection is therefore contractible. The nerve of this
cover is \(C_n*C_n\), where \(C_n\) is the simplicial cycle with \(n\)
vertices: the horizontal and vertical intersection conditions are
independent. The finite subcomplex nerve theorem
~\cite[Theorem~10.6(i)]{bjorner} gives
\[
L^-\simeq C_n*C_n\simeq S^1*S^1\cong S^3.
\]
\end{proof}

\section{Coordinate restrictions}

We will compare \(K^-\) and \(L^-\) on coordinate semicircles. A cyclic
interval of length \(\ell\) means
\(\{c,c+1,\ldots,c+\ell\}\), with all entries reduced modulo \(n\).

\begin{lemma}\label{lem:semicircle}
Every simplex of \(K\) has at least one coordinate projection contained
in a cyclic interval of length \(k\).
\end{lemma}

\begin{proof}
Embed the grid in \(\mathbb R^2\oplus\mathbb R^2\) by
\[
v(a,b)=
(\cos(\pi a/k),\sin(\pi a/k),
 \cos(\pi b/k),\sin(\pi b/k)).
\]
For angular coordinate distances \(\alpha,\beta\in[0,\pi]\), the pair
inner product is
\[
\cos\alpha+\cos\beta
=2\cos\frac{\alpha+\beta}{2}\cos\frac{\alpha-\beta}{2}.
\]
It is nonnegative exactly when \(\alpha+\beta\leq\pi\).
The cases where the second factor is zero also have
\(\alpha+\beta=\pi\). Thus clique vertices have pairwise nonnegative
inner products.

Write them as \((x_i,y_i)\). If the convex hull of the \(x_i\) excludes
zero, choose its closest point \(p\) to zero. Minimality along the segment
from \(p\) to \(x_i\) gives \(p\cdot x_i\geq |p|^2>0\);
the first coordinates lie in an open semicircle.
Otherwise choose nonnegative \(\lambda_i\), summing to one, with
\(\sum_i\lambda_i x_i=0\), and put \(w=\sum_i\lambda_i y_i\).
Since \(|(x_i,y_i)|^2=2\),
\[
|w|^2=\Big|\sum_i\lambda_i(x_i,y_i)\Big|^2
\geq 2\sum_i\lambda_i^2>0.
\]
Moreover \(y_j\cdot w\geq0\) for every \(j\), by the pair inner products.
The second coordinates consequently lie in a closed semicircle.
For integer coordinates, either semicircle can be enlarged to one of
the stated intervals.
\end{proof}

\subsection*{Finite cylinders}

For \(1\leq m\leq k+1\), let \(F_m\) be the closed Rips complex at radius
\(k\) on
\(\{0,\ldots,m-1\}\times\ZZ/n\ZZ\), with metric
\[
|x-x'|+\delta_n(y,y').
\]
Only the first coordinate is unwrapped.

\begin{proposition}\label{prop:cylinder}
For \(1\leq m\leq k+1\),
\[
F_m\simeq\bigvee^{m-1}S^{n-1},
\]
where a zero-fold wedge is a point. Deleting the selected facets of its
\(m-1\) adjacent-column bands leaves a contractible complex \(F_m^-\).
\end{proposition}

\begin{proof}
The case \(m=1\) is a simplex. For \(m\geq2\), put \(t=m-1\), and let
\[
U=\bigcup_{c\in\ZZ/n\ZZ}\st_{F_m}(t,c),
\qquad E=F_{m-1}\cap U.
\]
Here the closed star consists of faces that can be extended by the indicated
vertex. The vertices of the last column form a simplex. All nonempty
intersections of their closed stars are cones at any of the corresponding
last-column vertices. Their nerve is a full simplex, so \(U\) is
contractible, and \(F_m=F_{m-1}\cup U\).

Cover \(E\) by \(E_c=F_{m-1}\cap\st_{F_m}(t,c)\).
For each nonempty proper set \(S\subset\ZZ/n\ZZ\), there is a \(z\) with
\[
z\notin S+k,\qquad \delta_n(z,s_0)\leq1
\quad\text{for some }s_0\in S.                 \tag{*}
\]
If some \(s\in S\) has \(s+k\notin S\), take \(z=s\).
Otherwise \(S=S+k\), and a cyclic edge from \(S\) to its complement
gives \(s_0,z\) as required.

The vertex \((t-1,z)\) belongs to every \(E_c\), \(c\in S\), since
\(\delta_n(z,c)\leq k-1\). If \((x,y)\) belongs to their intersection,
put \(r=t-x\). Then \(r+\delta_n(y,s_0)\leq k\), and
\[
d((t-1,z),(x,y))
\leq r-1+\delta_n(z,s_0)+\delta_n(s_0,y)\leq k.
\]
Thus the intersection is a cone at \((t-1,z)\).
The intersection over all \(n\) centers is empty: the second-coordinate
antipode of any previous-column vertex gives distance \(r+k>k\).
The nerve is therefore the boundary of an \((n-1)\)-simplex, and
\(E\simeq S^{n-2}\).

Inductively \(F_{m-1}\) is \((n-2)\)-connected. This connectivity follows
from its wedge description and cellular approximation
~\cite[Theorem~4.8]{hatcher}. Both maps from \(E\) to \(F_{m-1}\) and
\(U\) are null-homotopic. The null-homotopic union lemma gives
\[
F_m\simeq F_{m-1}\vee\Sigma E
\simeq\bigvee^{m-1}S^{n-1}.
\]

The facet-removal argument of Lemma~\ref{lem:split} also applies inside
this cylinder:
\[
F_m\simeq F_m^-\vee\bigvee^{m-1}S^{n-1}.
\]
It follows that \(F_m^-\) has trivial reduced integral homology.
Its fundamental group is unchanged by the facet deletions, which occur
in dimension \(n-1\geq5\); cellular approximation applies to the common
two-skeleton. Hence \(F_m^-\) is simply connected.
The homology form of Whitehead's theorem
~\cite[Corollary~4.33]{hatcher}, applied to its map to a point, proves
contractibility.
\end{proof}

\subsection*{Rectangles and restricted bands}

\begin{lemma}\label{lem:rectangle}
The closed Manhattan Rips complex at radius \(k\) on a nonempty finite
integer rectangle is contractible.
\end{lemma}

\begin{proof}
By~\cite[Corollary~4.3]{torus}, this complex is the nerve of the open
Manhattan balls of radius \((k+1)/2\) centered at its grid points.
The balls have convex intersections and cover the whole rectangular region:
each point is within Manhattan distance one of a grid point.
Coordinate clamping to the rectangle decreases distance to every grid
center. The straight-line homotopy to the clamped point stays in each ball
containing the starting point. Thus the union deforms onto the rectangle,
and the nerve theorem~\cite[Theorem~10.7]{bjorner} proves the claim.
\end{proof}

For coordinate sets \(I,J\), write \(X|_{I\times J}\) for the induced
restriction of a complex \(X\).

\begin{lemma}\label{lem:restricted-bands}
If \(I,J\) are each a nonempty cyclic interval of length at most \(k\)
or the whole circle, and at least one is a proper interval, then
\(L^-|_{I\times J}\) is contractible.
\end{lemma}

\begin{proof}
Cover it by the restrictions of the punctured bands.
A nonempty band restricted to a proper interval in its circular
coordinate is a cone. In its
cross-polytope description, at least one vertex has its opposite absent:
an interval of length less than \(k\) has no antipodal pair, while one
of length \(k\) has interior vertices with antipodes outside.
No selected facet lies in this restricted direction, since its parity
projection has minimum cyclic span \(n-2>k\).
If a full circular direction is retained, a complete band in that
direction is its punctured ball. A restriction to one row or column is
a simplex. All nonempty intersections of distinct restricted bands remain
simplices.

For a full coordinate the band-intersection graph is an \(n\)-cycle.
For a proper interval it is a path: its vertices are the consecutive bands
meeting that interval, and its edges correspond to their shared
coordinates in the interval. This includes the two bands meeting a
singleton. The cover nerve is the join of the two graphs. It has a path
factor, so it is contractible. The subcomplex nerve theorem completes
the proof.
\end{proof}

\section{Comparison and the homotopy formula}

We use the following finite-cover homology observation.

\begin{lemma}\label{lem:comparison}
Suppose a finite simplicial complex \(X\) is covered by subcomplexes
\(X_1,\ldots,X_s\), and \(Y\subset X\). If for every nonempty index set
\(J\subset\{1,\ldots,s\}\) the inclusion
\[
Y\cap\bigcap_{j\in J}X_j\ \longrightarrow\ \bigcap_{j\in J}X_j
\]
induces an isomorphism on integral homology, then \(Y\hookrightarrow X\)
does so.
\end{lemma}

\begin{proof}
Induct on \(s\). Separate \(X_s\) from the union of the preceding
members. The induction hypothesis applies both to that union and to
its intersection with \(X_s\), covered by \(X_i\cap X_s\), \(i<s\).
Naturality of the Mayer--Vietoris sequences for these subcomplex unions
and the five lemma give the assertion in every degree; see
~\cite[Section~2.2, pp.~149--150]{hatcher}.
\end{proof}

\begin{proof}[Proof of Theorem~\ref{thm:main}]
By Lemma~\ref{lem:semicircle}, the induced restrictions to
\[
\{a,\ldots,a+k\}\times\ZZ/n\ZZ,\qquad
\ZZ/n\ZZ\times\{b,\ldots,b+k\},
\]
cover \(K^-\) simplex by simplex. Use the same restrictions for \(L^-\).

Every intersection is a restriction to \(I\times J\). Each coordinate
set is the whole circle if unconstrained; otherwise it is empty,
a connected cyclic interval of length at most \(k\), or a pair of
antipodal points. Indeed, two closed semicircles either meet in an
interval or, when complementary, in their two endpoints. Intersecting
a shorter interval with another semicircle cannot disconnect it, since
their total lengths are less than a full circumference. Further cuts of
an endpoint pair give a singleton or the empty set.

If one coordinate is an antipodal pair, every clique is contained in
one of its two rows or columns, or is a pair whose other coordinate is
constant. All these faces already lie in \(L\), and none is deleted.
Thus the restricted inclusion from \(L^-\) to \(K^-\) is the identity.
The same is true for a singleton coordinate.

If both coordinates are proper connected intervals, the metric unwraps
to an ordinary rectangle. No selected facet lies there because
\(n-2>k\). Lemmas~\ref{lem:rectangle} and~\ref{lem:restricted-bands}
make both restrictions contractible.
If exactly one coordinate is proper and the other is full,
Proposition~\ref{prop:cylinder} makes the restriction of \(K^-\)
contractible, and Lemma~\ref{lem:restricted-bands} does the same for
\(L^-\). Empty intersections occur on both sides.

Lemma~\ref{lem:comparison} now gives
\[
H_*(L^-;\ZZ)\cong H_*(K^-;\ZZ).
\]
By Proposition~\ref{prop:band-nerve}, the only nonzero reduced group of
\(K^-\) is \(H_3(K^-;\ZZ)\cong\ZZ\).

The complex \(K^-\) is simply connected. Each Rips edge is homotopic,
relative to its endpoints, to a shortest grid path of length at most
\(k\). Its path vertices form a retained simplex, having at most
\(k+1<n\) vertices. Unit-square triangulations are also retained.
Square homotopies commute horizontal and vertical steps, so a closed grid
loop reduces to powers of the two coordinate cycles. These cycles bound
in the retained full row and full column simplices.

Hurewicz first gives \(\pi_2(K^-)=0\) and then represents a generator of
\(H_3(K^-;\ZZ)\) by a map \(S^3\to K^-\).
This map is an integral homology isomorphism, so
~\cite[Theorem~4.32 and Corollary~4.33]{hatcher} gives \(K^-\simeq S^3\).
Lemma~\ref{lem:split} restores the \(2n=4k\) spheres of dimension
\(n-1=2k-1\), proving the theorem.
\end{proof}

\section*{AI assistance}
OpenAI Codex was used to explore and draft the proof, check the cited
theorem statements, and prepare the manuscript. Independent human
mathematical review is requested.

\begin{thebibliography}{9}
\bibitem{torus}
H.~Adams, A.~Y.~Adetowubo, H.~Barriga-Acosta, Z.~Feng and J.~Sterling,
\emph{Vietoris--Rips complexes of torus grids},
Mediterranean Journal of Mathematics \textbf{22} (2025), article 173,
\href{https://doi.org/10.1007/s00009-025-02945-9}{doi:10.1007/s00009-025-02945-9}.
Author version: \href{https://arxiv.org/abs/2502.07134v2}{arXiv:2502.07134v2}.
\bibitem{bjorner}
A.~Bj\"orner, \emph{Topological methods},
in R.~L.~Graham, M.~Gr\"otschel and L.~Lov\'asz (eds.),
\emph{Handbook of Combinatorics}, vol.~2, Elsevier, 1995, pp.~1819--1872.
\bibitem{hatcher}
A.~Hatcher, \emph{Algebraic Topology}, Cambridge University Press, 2002.
\end{thebibliography}
\end{document}
